\section{Discussion on Adaptive Speculative Strategies Derived from Empirical Findings}
\label{sec:acsd}

The current empirical study identifies the strongest fixed-policy speedup at
smaller block size ($k{=}4$), with net speed degrading as $k$ increases even
when accepted block length rises. This indicates that fixed-policy drafting is
vulnerable to acceptance/overhead imbalance and long-tail latency. This
motivates a derived adaptive strategy guideline:
Adaptive Cascaded Speculative Decoding (ACSD), which keeps the standard
verifier but adapts policy online. In the latest implementation, ACSD is
integrated into the Phase 7/8 pipeline with additional runtime-stability
controls to reduce pathological latency tails in long runs.

ACSD is inspired by a practical deployment gap. Advanced speculative variants
such as EAGLE-3 and DFlash reduce drafting cost through new drafter
architectures and dedicated training/integration pipelines, and are typically
reported on very strong accelerators. Our objective is different: retain a
same-family drafter/verifier stack that can run on consumer hardware, then
recover additional speed by adapting runtime policy (adaptive $k$, rescue
switching, and tail fallback) instead of introducing a new trained drafter.

\subsection{Design objective}

Let speculative-cycle latency be
\begin{equation}
  L = \frac{T_{\text{draft}} + T_{\text{verify}}}{\tau},
\end{equation}
where $\tau$ is expected accepted tokens per cycle. In fixed-policy speculative
decoding, poorly matched samples can reduce $\tau$ and increase total verify
steps, producing heavy-tail runtime. ACSD explicitly targets this effect by
adapting draft policy at sample level.

\subsection{Proposed approach}

ACSD combines four online controls while preserving standard target-side
verification:
\begin{enumerate}
      \item \textbf{Adaptive $k$:} choose $k\in\{3,4,5\}$ from rolling acceptance,
        using larger blocks only when acceptance supports it. To avoid cold-start
        overshoot, each sample starts from base $k$ before adaptive updates.
  \item \textbf{Cascaded draft switch:} use 0.5B as a fast primary drafter and
        switch to 1.5B rescue drafting when acceptance remains low for
        consecutive verify steps, then enforce a cooldown before a new rescue
        episode to reduce oscillation.
  \item \textbf{Tail guardrail:} if acceptance stays below a threshold after a
        minimum generated length for multiple consecutive checks, fall back to
        target-only autoregressive decoding for the remainder of that sample.
  \item \textbf{Checkpointed execution:} persist partial CSV and verify-log
        artifacts every fixed number of samples so interrupted long runs retain
        recoverable progress.
\end{enumerate}

This design preserves verifier compatibility and allows direct comparison with
vanilla speculative decoding under the same manifests, prompts, metrics, and
hardware.

\begin{figure}[t]
\centering
\resizebox{0.98\columnwidth}{!}{%
\begin{tikzpicture}[
      node distance=4.5mm and 5mm,
      box/.style={rectangle, draw, rounded corners=1mm, align=center, minimum height=6mm, text width=30mm, inner sep=1.5mm},
      decision/.style={diamond, draw, align=center, aspect=2, inner sep=1.1mm, text width=22mm},
      line/.style={-Latex, thick, shorten >=1pt, shorten <=1pt}
]
\node[box] (start) {Start sample\\init primary draft, base $k$};
\node[box, below=of start] (draft) {Draft proposes $k$ tokens\\(0.5B primary or 1.5B rescue)};
\node[box, below=of draft] (verify) {Target verifies draft block\\compute rolling $\alpha$};
\node[decision, below=of verify] (high) {$\alpha \ge \tau_{\text{high}}$?};
\node[box, left=of high, xshift=-8mm] (kup) {Increase $k$\\toward 5};
\node[decision, below=of high] (low) {$\alpha \le \tau_{\text{low}}$?};
\node[box, left=of low, xshift=-8mm] (kdown) {Reduce $k$\\toward 3};
\node[decision, below=of low] (rescue) {Low-accept streak\\$\ge c$ and cooldown=0?};
\node[box, left=of rescue, xshift=-8mm] (switch) {Switch to 1.5B rescue\\for $h$ cycles};
\node[decision, below=of rescue] (fallback) {$\alpha < \tau_{\text{AR}}$ for\\$m$ checks and length $\ge L_{\min}$?};
\node[box, right=of fallback, xshift=8mm] (ar) {Fallback to target-only\\AR for tail};
\node[box, below=of fallback] (loop) {Continue next ACSD cycle};
\node[box, below=of loop] (end) {EOS or max tokens\\end sample};
\coordinate (loopRightTop) at ([xshift=18mm]draft.east);
\coordinate (loopRightBottom) at ([xshift=18mm]end.east);

\draw[line] (start) -- (draft);
\draw[line] (draft) -- (verify);
\draw[line] (verify) -- (high);
\draw[line] (high) -- node[above, sloped, font=\scriptsize] {yes} (kup);
\draw[line] (kup.east) to[out=0, in=180] (draft.west);
\draw[line] (high) -- node[right, font=\scriptsize] {no} (low);
\draw[line] (low) -- node[above, sloped, font=\scriptsize] {yes} (kdown);
\draw[line] (kdown.east) to[out=0, in=180] (draft.west);
\draw[line] (low) -- node[right, font=\scriptsize] {no} (rescue);
\draw[line] (rescue) -- node[above, sloped, font=\scriptsize] {yes} (switch);
\draw[line] (switch.east) to[out=0, in=180] (draft.west);
\draw[line] (rescue) -- node[right, font=\scriptsize] {no} (fallback);
\draw[line] (fallback) -- node[above, font=\scriptsize] {yes} (ar);
\draw[line] (fallback) -- node[right, font=\scriptsize] {no} (loop);
\draw[line] (ar) |- (end.east);
\draw[line] (loop) -- (end);
\draw[line] (end.east) -- (loopRightBottom) -- (loopRightTop) -- (draft.east);
\end{tikzpicture}%
}
\caption{Adaptive Cascaded Speculative Decoding (ACSD) control loop. ACSD
adapts block size from rolling acceptance, switches from the 0.5B primary to a
1.5B rescue drafter when acceptance remains low, applies rescue cooldown, and
uses delayed target-only autoregressive fallback for pathological tails.}
\label{fig:acsd-flow}
\end{figure}

\subsection{Evaluation plan}

The extension is evaluated as a strict incremental study rather than a new
protocol. We retain:
\begin{itemize}
  \item the same target model and task manifests,
  \item the same deterministic and stochastic regimes,
  \item the same reporting metrics ($S$, $\alpha$, $B_{\text{eff}}$, quality deltas),
  \item the same artifact schema for per-configuration CSV outputs, with ACSD
        metadata fields for adaptive-$k$ usage, rescue usage, and fallback rate.
\end{itemize}

Primary comparison points are the current best fixed-policy speculative
setting ($k{=}4$ deterministic) and the high-$k$ settings ($k{=}16$) where
fixed-policy decoding showed diminishing returns.

\subsection{Empirical ACSD results}

We ran ACSD on the full 1{,}000-sample manifests for both deterministic and
stochastic regimes, using the same target model and evaluation protocol as the
fixed-policy grid. Table~\ref{tab:acsd-headline} reports headline runtime
metrics. Pairwise speedup $S$ is recomputed from per-sample baseline and ACSD
latencies matched by sample identifier and task.

\begin{table}[t]
\centering
\caption{ACSD headline results on RTX4090 (0.5B primary, 1.5B rescue,
adaptive $k\in\{3,4,5\}$). $\Delta S_{\text{best}}$ is ACSD speedup minus the
best fixed-policy speedup in the same regime.}
\label{tab:acsd-headline}
\small
\begin{tabular}{l c c c c c c c}
\hline
\textbf{Regime} & $S$ & $\alpha$ & $B_{\text{eff}}$ & $R_{\text{tok}}$ & Rescue & Fallback & $\Delta S_{\text{best}}$ \\
\hline
deterministic & 2.9792 & 0.4207 & 1.7977 & 32.61 & 0.698 & 1.6\% & -0.0882 \\
stochastic    & 2.8981 & 0.4071 & 1.7214 & 31.08 & 0.843 & 2.2\% & -0.0413 \\
\hline
\end{tabular}
\end{table}

The ACSD runtime is close to the best fixed-policy settings: $2.9\%$ below
best fixed deterministic speedup ($3.0674$) and $1.4\%$ below best fixed
stochastic speedup ($2.9394$). Rescue activation remains limited
(0.698/0.843 mean rescue steps per sample) and tail fallback is rare
(1.6\%/2.2\%), indicating that adaptive controls are used as guardrails rather
than dominating the decode path.

Quality remains near baseline in deterministic ACSD
($\Delta$GSM8K $=+0.33$, $\Delta$MMLU $=0.00$, $\Delta$CNN/DM $=-0.02$), while
stochastic ACSD shows moderate degradation
($\Delta$GSM8K $=-1.00$, $\Delta$MMLU $=-0.20$, $\Delta$CNN/DM $=-0.22$).
Aggregated adaptive-$k$ traces show most cycles at $k\in\{4,5\}$
(deterministic: $86.9\%$; stochastic: $85.8\%$), with $k{=}3$ used as a
recovery mode and very short steps ($k\in\{1,2\}$) mainly near completion.
